% !TeX program = xelatex
\documentclass[11pt,a4paper]{article}
\usepackage[margin=25mm,headheight=14pt]{geometry}
\usepackage{fontspec}
\setmainfont{texgyretermes}[Extension=.otf,UprightFont=*-regular,
 BoldFont=*-bold,ItalicFont=*-italic,BoldItalicFont=*-bolditalic]
\setmonofont{lmmono10-regular.otf}[Scale=0.86]
\usepackage{booktabs,tabularx,longtable,array}
\usepackage{amsmath,amssymb}
\usepackage{enumitem,fancyhdr,titlesec}
\usepackage[unicode,colorlinks=true,linkcolor=black,
 urlcolor=black,citecolor=black]{hyperref}
\usepackage{xurl}

\titleformat{\section}{\large\bfseries}{\thesection}{0.7em}{}
\titleformat{\subsection}{\normalsize\bfseries}{\thesubsection}{0.7em}{}
\titlespacing*{\section}{0pt}{14pt}{6pt}
\titlespacing*{\subsection}{0pt}{10pt}{4pt}

\setlength{\parindent}{0pt}
\setlength{\parskip}{5pt}
\setlength{\emergencystretch}{2em}
\linespread{1.06}
\setlist{leftmargin=1.6em,itemsep=2pt,topsep=3pt}
\renewcommand{\arraystretch}{1.15}
\renewcommand{\contentsname}{Mục lục}
\setcounter{tocdepth}{1}
\renewcommand{\refname}{Tài liệu tham khảo}

\newcommand{\pathref}[1]{\nolinkurl{#1}}
\newcommand{\AMO}{\operatorname{AMO}}
\newcommand{\LB}{\mathrm{LB}}
\newcommand{\UB}{\mathrm{UB}}
\newcommand{\file}[1]{\texttt{#1}}

\pagestyle{fancy}
\fancyhf{}
\fancyhead[L]{\small Hướng dẫn nâng cấp MaxSAT-DDD}
\fancyhead[R]{\small Bản dành cho sinh viên --- v3.0}
\fancyfoot[C]{\thepage}
\hypersetup{pdftitle={Hướng dẫn sinh viên nâng cấp MaxSAT-DDD},
 pdfauthor={Tài liệu hướng dẫn sinh viên}}

% =============================================================================
\begin{document}
\hypersetup{pageanchor=false}
\begin{titlepage}
\vspace*{20mm}
{\large TÀI LIỆU HƯỚNG DẪN SINH VIÊN\par}
\vspace{12mm}
{\LARGE\bfseries Nâng cấp MaxSAT-DDD\par}
\vspace{4mm}
{\Large cho bài toán điều độ lại tàu\par}
\vspace{10mm}
{\large Hướng dẫn triển khai theo từng giai đoạn\par}
\vspace{10mm}
\rule{\textwidth}{0.4pt}
\vspace{4mm}

Tài liệu này hướng dẫn sinh viên mở rộng kết quả của bài hội nghị~\cite{conference}
theo hướng nộp tạp chí \emph{Computers \& Operations Research}.
Tài liệu được tổ chức theo đúng thứ tự thực hiện:
\begin{enumerate}
\item Đọc và chạy được code hiện tại (Mục~\ref{sec:read}).
\item Hiểu bài toán và ký hiệu (Mục~\ref{sec:problem}).
\item Chốt kết quả nền để so sánh (Mục~\ref{sec:baseline}).
\item Cài đặt từng cải tiến và kiểm thử (Mục~\ref{sec:improve1}--\ref{sec:improve3}).
\item Tích hợp, thực nghiệm và viết báo cáo (Mục~\ref{sec:integrate}--\ref{sec:report}).
\end{enumerate}

Mỗi phần đều bắt đầu bằng ví dụ số cụ thể \emph{trước} khi đưa ra công thức.
Sinh viên nên đọc tuần tự, không bỏ qua bước kiểm thử.

\vfill
{\normalsize Phiên bản 3.0 --- tháng 9 năm 2026\par}
\vspace{3mm}
Nguồn nền: \pathref{manuscript/submission/main.tex} và mã nguồn tại \pathref{maxsattrainscheduling/}.
\end{titlepage}
\hypersetup{pageanchor=true}
\begingroup\small\setlength{\parskip}{0pt}\tableofcontents\endgroup
\newpage

% =============================================================================
\section{Đọc code trước khi làm bất cứ điều gì}
\label{sec:read}
% =============================================================================

Đọc 5 file sau theo đúng thứ tự. Không cần đọc toàn bộ mỗi file --- chỉ cần
hiểu phần được chỉ định.

\begin{tabularx}{\textwidth}{@{}p{9mm}p{62mm}X@{}}
\toprule
TT & File & Cần hiểu gì \\\midrule
1 & \file{src/problem.rs} &
  Struct \texttt{Visit}, \texttt{Train}, \texttt{Problem};
  hàm \texttt{visit\_delay\_cost()} tính chi phí trễ. \\
2 & \file{src/solvers/ddd/shared/precedence.rs} &
  Hàm \texttt{chain\_earliest()}: lan truyền thời gian dọc hành trình. \\
3 & \file{src/solvers/ddd/shared/common.rs} &
  Struct \texttt{Occ} và trường \texttt{delays}: nơi lưu các mốc thời gian. \\
4 & \file{src/solvers/ddd/maxsat\_rc2.rs} &
  Vòng \texttt{loop\{...\}} từ dòng 139. Đây là file cần sửa nhiều nhất. \\
5 & \file{src/solvers/util/heuristic.rs} &
  Hàm \texttt{solve\_heuristic\_better()}: sinh lịch hợp lệ $\UB$. \\
\bottomrule
\end{tabularx}

\medskip
Câu hỏi cần trả lời được sau khi đọc:
\begin{itemize}
\item Dữ liệu đi vào từ đâu? Struct nào nhận?
\item Một mốc thời gian được thêm vào model ở dòng nào?
\item Chi phí trễ được tính và gửi vào solver ở đâu?
\item Chương trình biết khi nào thì dừng?
\end{itemize}

\textit{Kết quả cần có:} Chạy được ít nhất 1 instance nhỏ và in ra số vòng lặp,
số mốc thêm mỗi vòng, chi phí cuối cùng.

% =============================================================================
\section{Hiểu bài toán}
\label{sec:problem}
% =============================================================================

\subsection{Bài toán là gì?}

Ba tàu bị trễ do sự cố. Mỗi tàu phải đi qua một chuỗi đoạn đường và ga.
Một số đoạn đường chỉ cho 1 tàu đi cùng lúc (đường đơn). Chương trình cần
quyết định thứ tự và thời điểm sao cho: (1) không có 2 tàu nào chiếm cùng
đoạn đường đơn, và (2) tổng thời gian trễ là nhỏ nhất.

\subsection{DDD hoạt động như thế nào?}

Thay vì xét mọi giây có thể, DDD bắt đầu với rất ít mốc thời gian rồi
tăng dần theo vòng lặp:

\begin{enumerate}
\item Tạo mô hình nhỏ (ít mốc thời gian).
\item Giải MaxSAT $\to$ nhận lịch ứng viên.
\item Kiểm tra lịch: nếu không còn vi phạm $\to$ dừng, trả nghiệm tối ưu.
\item Thêm mốc thời gian ở chỗ vi phạm $\to$ quay lại bước 1.
\end{enumerate}

Bước 4 gọi là \textbf{refinement} --- đây là phần đề tài muốn cải thiện.

\subsection{Ba loại hàm chi phí trong code}

Xem \file{src/problem.rs}, enum \texttt{DelayCostType}.

\begin{tabularx}{\textwidth}{@{}p{38mm}p{30mm}X@{}}
\toprule
Tên trong code & Cách tính phạt & Ví dụ (trễ 200s, $\Delta{=}180$s) \\\midrule
\texttt{FiniteSteps*} & Phạt cố định khi vượt ngưỡng
  & Trễ 200s: phạt 1; trễ 10000s: vẫn 1 \\
\texttt{InfiniteSteps180} & Mỗi 180s trễ tăng thêm 1 bậc
  & Trễ 200s: phạt $\lceil200/180\rceil=2$ \\
\texttt{Continuous} & Mỗi giây trễ = 1 điểm phạt
  & Trễ 200s: phạt 200 \\
\bottomrule
\end{tabularx}

\medskip
Lưu ý quan trọng: với \texttt{FiniteSteps}, phạt bị chặn trên nên không thể
suy ra giới hạn thời gian từ ngân sách chi phí. Với hai loại còn lại thì có thể.
Điều này sẽ quan trọng ở Mục~\ref{sec:cost-bound}.

\subsection{Ký hiệu dùng trong các công thức}
\label{sec:notation}

Mỗi hoạt động được đánh dấu bằng mã $v = (\text{tàu},\,\text{vị trí trên hành trình})$.

\begin{tabularx}{\textwidth}{@{}p{28mm}X@{}}
\toprule
Ký hiệu & Ý nghĩa \\\midrule
$t_v$ & Thời điểm bắt đầu hoạt động $v$ (biến cần tìm). \\
$e_v = t_{s(v)}$ & Thời điểm rời tài nguyên = thời điểm hoạt động kế tiếp bắt đầu. \\
$L_v,\;U_v$ & Sớm nhất và muộn nhất $t_v$ được phép ($L_v \le t_v \le U_v$). \\
$p_v$ & Thời gian chạy tối thiểu (travel time). \\
$s(v)$ & Hoạt động kế tiếp cùng tàu. \\
$a_v$ & Thời điểm theo lịch ban đầu (aimed time). \\
$d_v(\tau)$ & Biến Boolean: đúng khi $t_v \ge \tau$. \\
\bottomrule
\end{tabularx}

\medskip
Ba ràng buộc cơ bản:
\begin{align}
 t_{s(v)} &\ge t_v + p_v
   &&\text{(tàu phải ở đủ thời gian trên tài nguyên)}\label{eq:travel}\\
 t_w \ge e_v \;\;&\text{hoặc}\;\; t_v \ge e_w
   &&\text{(hai tàu không chiếm cùng tài nguyên)}\label{eq:resource}\\
 F(t) &= \sum_{v \in M} c_v\!\bigl(\max(0,\,t_v - a_v)\bigr)
   &&\text{(tổng chi phí trễ cần tối thiểu hóa)}\label{eq:objective}
\end{align}

% =============================================================================
\section{Giai đoạn 0: chốt kết quả nền}
\label{sec:baseline}
% =============================================================================

Trước khi thêm bất kỳ cải tiến nào, cần xác định chính xác phiên bản code và
dữ liệu đã tạo ra kết quả của bài hội nghị.

\begin{tabularx}{\textwidth}{@{}p{42mm}X@{}}
\toprule
Điểm cần kiểm tra & Việc cần làm \\\midrule
Số bài CPLEX giải được & Đối chiếu số trong bảng bài báo và file CSV. \\
Mã hóa SC-AMO & Xác định đúng phiên bản code và ngưỡng đã dùng. \\
Đơn vị thời gian & Kiểm tra 180 là giây hay phút trong dữ liệu và code. \\
Ý nghĩa track/station & Dữ liệu thay đổi gì: travel time, dwell time, hay cấu trúc tài nguyên? \\
Giới hạn chạy & Ghi rõ timeout, giới hạn RAM, số luồng, phiên bản solver. \\
\bottomrule
\end{tabularx}

\medskip
Chạy script kiểm tra (chỉ đọc, không sửa):
\begin{center}
\texttt{python3 manuscript/journal/audit\_existing\_results.py}
\end{center}

\textit{Kết quả cần có:} Một cấu hình nền chạy lại được; file ghi rõ cấu hình;
bảng ánh xạ tên instance giữa CSV và code.

% =============================================================================
\section{Cải tiến 1: suy luận sớm về thời gian}
\label{sec:improve1}
% =============================================================================

\subsection{Vấn đề trong code hiện tại}

Tại \file{maxsat\_rc2.rs} dòng 96--100, mỗi hoạt động được khởi tạo với:

\begin{verbatim}
delays: vec![(true,  visit.earliest),  // L_v = gio som nhat
             (false, i32::MAX)],       // U_v = VO CUC
\end{verbatim}

\texttt{i32::MAX} $\approx$ 68 năm. Chương trình không biết tàu muộn nhất
được phép đến lúc nào, nên DDD có thể thêm mốc ở những thời điểm vô nghĩa.

Thêm vào đó, giá trị $\UB$ (chi phí lịch tốt nhất) ở dòng 141 chỉ dùng để
in khi timeout, chưa được dùng để siết $U_v$ cho từng hoạt động.

\subsection{Ví dụ số trước khi xem công thức}

\textbf{Ví dụ 1: Lan truyền trong cùng một tàu.}

Tàu A có 3 điểm dừng với travel time 5s và 3s. Thời điểm sớm nhất ban đầu:
$L_0 = 10$, $L_1 = 8$, $L_2 = 20$.

Sau lan truyền forward:
\[
 L_1 \leftarrow \max(8,\; 10+5) = 15,\qquad
 L_2 \leftarrow \max(20,\; 15+3) = 20.
\]

Tàu A không thể đến điểm 1 trước giây 15, dù \texttt{visit.earliest} ghi là 8.

\medskip
\textbf{Ví dụ 2: Suy ra thứ tự bắt buộc giữa 2 tàu.}

Tàu A và tàu B cùng cần đoạn đường X:
\begin{itemize}
\item $e_A^{\min} = 15$: tàu A thoát đường X sớm nhất lúc giây 15.
\item $U_B = 13$: tàu B muộn nhất được vào lúc giây 13.
\end{itemize}

Nếu A đi trước B: B phải vào sau giây 15, nhưng $U_B = 13 < 15$. Không thể!
Vậy B \emph{bắt buộc} đi trước A. Thêm ngay ràng buộc này vào solver.

\subsection{Mức 1: lan truyền $L_v$ và $U_v$}
\label{sec:prop}

Hàm cần viết: \file{shared/precedence.rs}, hàm \texttt{propagate\_bounds()}.
(Hàm \texttt{chain\_earliest()} đã có sẵn; mở rộng từ đó.)

Với bất kỳ ràng buộc $t_b \ge t_a + h$ nào, cập nhật hai chiều:
\begin{align}
 L_b &\leftarrow \max(L_b,\; L_a + h), \label{eq:prop-lower}\\
 U_a &\leftarrow \min(U_a,\; U_b - h). \label{eq:prop-upper}
\end{align}

Lặp đến khi không có giá trị nào thay đổi nữa. Cho cặp liên tiếp trong cùng
một tàu: dùng $a = v$, $b = s(v)$, $h = p_v$.

\textbf{Ví dụ tính tay:} $L_a=10$, $L_b=8$, $U_a=30$, $U_b=20$, $h=5$:
\[
 L_b = \max(8,\; 10+5) = 15,\qquad
 U_a = \min(30,\; 20-5) = 15.
\]

\textbf{Cài đặt trong \file{maxsat\_rc2.rs}:}
\begin{itemize}
\item Sau dòng 112 (sau vòng khởi tạo occupations): khởi tạo \texttt{eff\_lb}
      bằng \texttt{chain\_earliest()}, \texttt{eff\_ub} bằng \texttt{i32::MAX};
      gọi \texttt{propagate\_bounds()} lần đầu.
\item Mỗi vòng lặp, sau khi phát hiện vi phạm và sau khi $\UB$ được cập nhật:
      gọi lại \texttt{propagate\_bounds()}.
\end{itemize}

Lưu ý: dùng \texttt{saturating\_add} / \texttt{saturating\_sub} thay vì
\texttt{+} / \texttt{-} khi làm việc với \texttt{i32::MAX} để tránh tràn số.

\subsection{Tìm thứ tự bắt buộc giữa hai tàu}
\label{sec:forced}

Hàm cần viết: \file{shared/precedence.rs}, hàm \texttt{find\_forced\_orders()}.

Gọi $e_v^{\min} = L_{s(v)}$ là thời điểm sớm nhất tàu rời tài nguyên.
Với hai hoạt động $v$, $w$ cùng xung đột (cùng cần 1 đoạn đường):
\begin{align}
 e_v^{\min} > U_w
 &\;\Longrightarrow\; \text{loại thứ tự } v \text{ trước } w,\label{eq:reject-vw}\\
 e_w^{\min} > U_v
 &\;\Longrightarrow\; \text{loại thứ tự } w \text{ trước } v.\label{eq:reject-wv}
\end{align}

Dùng dấu $>$ không phải $\ge$: hai tàu vừa khít vẫn là hợp lệ.

\begin{itemize}
\item Chỉ loại được~\eqref{eq:reject-vw}: $w$ bắt buộc đi trước $v$.
\item Chỉ loại được~\eqref{eq:reject-wv}: $v$ bắt buộc đi trước $w$.
\item Loại cả hai: bài toán vô nghiệm trong phạm vi đang xét.
\item Không loại được gì: giữ cả hai, chưa kết luận.
\end{itemize}

\textbf{Thêm ràng buộc cứng vào solver.} Xem hàm \texttt{Occ::time\_point()}
tại \file{shared/common.rs} dòng 54--94:

\begin{verbatim}
// Buoc t_v2 >= tau_exit
let (var, _) = occupations[v2_id].time_point(&mut solver, tau_exit);
solver.add_clause(vec![var]);  // unit clause = bat buoc TRUE
\end{verbatim}

\subsection{Mức 2: siết $U_v$ từ ngân sách chi phí}
\label{sec:cost-bound}

Hàm cần viết: \file{shared/precedence.rs}, hàm \texttt{tighten\_ub\_from\_cost()}.

\textbf{Chỉ làm sau khi Mức 1 đã đúng và ổn định.
Chỉ áp dụng với \texttt{Continuous} và \texttt{InfiniteSteps}.}

Ý tưởng: nếu heuristic tìm được lịch chi phí $B$, thì nghiệm tối ưu không
thể đắt hơn $B$. Từ đó ta có thể suy ra tàu $v$ được phép trễ tối đa bao lâu.

\begin{enumerate}
\item Phạt tối thiểu của mọi hoạt động $u \ne v$:
      $\ell_u = c_u\!\bigl(\max(0,\; L_u - a_u)\bigr)$.
\item Ngân sách còn lại cho $v$:
      $R_v = B - \sum_{u \ne v} \ell_u$.
\item Trễ tối đa cho phép:
      $D_v^{\mathrm{Cont}} = \lfloor R_v/w_v \rfloor$,\quad
      $D_v^{\mathrm{Round}} = \Delta\lfloor R_v/w_v \rfloor$.
\item Cập nhật: $U_v \leftarrow \min(U_v,\; a_v + D_v)$.
\end{enumerate}

\textbf{Ví dụ tính tay.} $B=50$, 2 tàu, Continuous.
Tàu 0: aimed=100, $L_0=100$ nên $\ell_0=0$.
Tàu 1: aimed=200, $L_1=210$ nên $\ell_1=10$.
\[
 R_0 = 50 - 10 = 40 \;\Rightarrow\; U_0 = 100 + 40 = 140.
 \qquad
 R_1 = 50 - 0 = 50 \;\Rightarrow\; U_1 = 200 + 50 = 250.
\]
Trước: $U_0 = U_1 = i32{::}MAX$. Sau: $U_0=140$, $U_1=250$.

\textbf{Cài đặt:} Sau khi \texttt{best\_heur} được cập nhật (dòng 190 của
\file{maxsat\_rc2.rs}), gọi \texttt{tighten\_ub\_from\_cost()}, rồi gọi lại
\texttt{propagate\_bounds()} để lan truyền $U_v$ mới ngược về phía trước.

\subsection{Các lỗi hay gặp}
\begin{itemize}
\item Dùng $\ge$ thay $>$ khi kiểm tra loại thứ tự --- trường hợp bằng nhau vẫn hợp lệ.
\item Không tính lại $U_v$ khi $B$ giảm (nghiệm heuristic tốt hơn xuất hiện).
\item Tràn số khi cộng \texttt{i32::MAX} mà không dùng \texttt{saturating\_add}.
\item Áp dụng Mức 2 cho \texttt{FiniteSteps} --- sẽ cho $U_v$ sai vì phạt không tăng theo thời gian.
\end{itemize}

\textit{Kết quả cần có:} Thành phần suy luận bật/tắt được; log số thứ tự bắt buộc
mỗi vòng; kết quả cuối cùng giống hệt baseline.

% =============================================================================
\section{Cải tiến 2: refinement có chọn lọc}
\label{sec:improve2}
% =============================================================================

\subsection{Vấn đề trong code hiện tại}

Tại \file{maxsat\_rc2.rs} dòng 574:

\begin{verbatim}
for (visit, var, t) in new_time_points.drain(..) {
    // them tat ca moc vi pham vao solver
}
\end{verbatim}

Sau khi phát hiện vi phạm (dòng 196--480), mọi mốc thời gian liên quan đều
được thêm ngay. Một vòng có 10 vi phạm có thể thêm 50+ mốc, khiến model
phình nhanh và solver chậm dần.

\subsection{Ví dụ số: tại sao cần chọn lọc}

Vòng 3 phát hiện 5 vấn đề. Có 3 gói ứng viên:

\begin{tabularx}{\textwidth}{@{}lXXX@{}}
\toprule
Gói & Giải quyết được & Mốc cần thêm & Điểm $= |C|/(1+m)$ \\\midrule
A & 3 vấn đề & 2 mốc & $3/3 = 1{,}0$ \\
B & 2 vấn đề & 3 mốc & $2/4 = 0{,}5$ \\
C & 1 vấn đề & 0 mốc mới & $1/1 = 1{,}0$ \\
\bottomrule
\end{tabularx}

\medskip
Chọn A và C: chỉ cần thêm 2 mốc. Vòng tiếp theo xử lý 2 vấn đề còn lại.
Kết quả cuối giống hệt khi thêm tất cả, nhưng mỗi vòng ít tốn kém hơn.

\subsection{Ba loại vi phạm cần nhận biết}

\begin{tabularx}{\textwidth}{@{}p{34mm}X@{}}
\toprule
Loại & Ý nghĩa \\\midrule
Hành trình &
  Một tàu đến hoạt động sau quá sớm so với hoạt động trước (vi phạm ràng buộc~\eqref{eq:travel}). \\
Tài nguyên &
  Hai tàu chiếm cùng tài nguyên trong cùng khoảng thời gian (vi phạm~\eqref{eq:resource}). \\
Chi phí &
  Lịch hợp lệ nhưng mô hình chưa biểu diễn đủ các mức chi phí thật.
  Không còn xung đột không có nghĩa là đã tối ưu. \\
\bottomrule
\end{tabularx}

\subsection{Tách phát hiện khỏi áp dụng}

Hiện tại phát hiện và thêm mốc xảy ra trong cùng một vòng lặp. Cần tách
thành 3 pha:

\begin{enumerate}
\item \textbf{Phát hiện:} duyệt tất cả hoạt động, ghi danh sách vi phạm.
      Không tạo biến mới, không sửa solver.
\item \textbf{Lập danh sách ứng viên:} với mỗi vi phạm, tạo một
      \texttt{RefinementPackage} chứa danh sách vấn đề nó xử lý và danh
      sách mốc cần thêm.
\item \textbf{Chọn lọc và áp dụng:} tính điểm, chọn $K$ mốc tốt nhất,
      thêm vào solver.
\end{enumerate}

\textbf{Mã giả tổng thể của vòng lặp DDD sau khi tách 3 pha:}

\begin{verbatim}
-- Moi vong lap DDD --

packages = []

-- PHA 1a: PHAT HIEN VI PHAM HANH TRINH (travel-time)
--   doc du lieu, KHONG goi time_point(), KHONG them clause
for each visit v in touched_intervals:
    t_in  = incumbent_time(v)
    t_out = incumbent_time(next(v))
    if t_in + travel_time(v) > t_out:
        t_in_var = delays[v][incumbent_idx].var   -- lit hien tai
        new_t    = t_in + travel_time(v)
        packages.add( TravelPackage(
            violation  = TravelTime(v),
            points     = [(next(v), new_t)],
            t1_in_var  = t_in_var,
            new_t      = new_t
        ))

-- PHA 1b: PHAT HIEN XUNG DOT TAI NGUYEN (resource conflict)
--   doc du lieu, KHONG goi time_point(), KHONG them clause
for each (v, w) in conflicting pairs:
    t1_out = incumbent_time(next(v))
    t2_out = incumbent_time(next(w))
    if intervals [t1_in, t1_out) and [t2_in, t2_out) overlap:
        t1_out_lit = delays[next(v)][incumbent_idx].var
        t2_out_lit = delays[next(w)][incumbent_idx].var
        packages.add( ResourcePackage(
            violation    = Resource(v, w),
            points       = [(w, t1_out), (v, t2_out)],
            t1_out_lit   = t1_out_lit,
            t2_out_lit   = t2_out_lit
        ))

-- PHA 2: CHON LOC (select_packages)
selected = greedy_select(packages, budget=K)
-- neu selected rong: buoc chon goi cu nhat (fallback)

-- PHA 3: AP DUNG goi da chon (apply_package)
for each package p in selected:
    if p is TravelPackage:
        var, is_new = time_point(next(v), p.new_t)  -- goi solver
        add_clause([NOT p.t1_in_var, var])           -- implication
        if is_new: new_time_points.add((next(v), var, p.new_t))
    if p is ResourcePackage:
        delay_w, _ = time_point(w, p.t1_out)        -- goi solver
        delay_v, _ = time_point(v, p.t2_out)        -- goi solver
        add_clause([NOT t1_out_lit, NOT t2_out_lit,  -- AMO clause
                    delay_v, delay_w])
        if is_new: new_time_points.add(...)
\end{verbatim}

\textbf{Tại sao tách pha?} Ở code cũ, phát hiện và thêm mốc xảy ra cùng lúc nên
không thể biết tổng chi phí của một tập vi phạm trước khi áp dụng. Sau khi tách,
pha 1 chỉ đọc; pha 2 quyết định; pha 3 mới ghi. Ba pha độc lập, có thể thay thuật toán chọn lọc mà không ảnh hưởng phần còn lại.

Struct cần tạo (thêm vào \file{shared/refinement.rs} mới):
\begin{verbatim}
pub enum ViolationId {
    TravelTime(VisitId),
    Resource(VisitId, VisitId),
    MissingCost(VisitId),
}
pub struct RefinementPackage {
    pub resolves:    Vec<ViolationId>,
    pub time_points: Vec<(VisitId, i32)>,
    pub first_seen:  usize,
}
\end{verbatim}

\subsection{Công thức tính điểm và chọn gói}

Gọi $C_{\mathrm{sel}}$ và $P_{\mathrm{sel}}$ là tập vấn đề và mốc đã được
chọn trong vòng hiện tại. Điểm của gói $b$:
\begin{equation}
 \operatorname{score}(b)
 = \frac{\bigl|C_b \setminus C_{\mathrm{sel}}\bigr|}
        {1 + \bigl|P_b \setminus P_{\mathrm{sel}}\bigr|}.
\label{eq:score}
\end{equation}

Tử số: số vấn đề mới gói xử lý được.
Mẫu số: số mốc mới cần thêm (cộng 1 tránh chia 0).

Vòng chọn lọc (lặp cho đến khi hết ngân sách $K$):
\begin{enumerate}
\item Tìm gói có điểm cao nhất và $\operatorname{score} > 0$.
\item Kiểm tra ngân sách: nếu tổng mốc sẽ vượt $K$ thì dừng.
\item Cập nhật $C_{\mathrm{sel}}$ và $P_{\mathrm{sel}}$.
\item Lặp lại.
\end{enumerate}

\textbf{Mã giả hàm \texttt{select\_packages}:}

\begin{verbatim}
fn select_packages(packages, budget K):
    if K = None: return all indices  -- budget vo han, them tat ca

    selected   = []      -- chi so goi da chon
    sel_viol   = {}      -- tap vi pham da duoc xu ly
    sel_points = {}      -- tap moc da duoc them
    total_pts  = 0       -- so moc moi da dung

    loop:
        best_score = -1
        best_idx   = None

        for each package p[i] not yet selected:
            if p[i].violation in sel_viol: skip  -- vi pham da co roi
            new_pts = count points in p[i].points NOT in sel_points
            if total_pts + new_pts > K: skip      -- vuot ngan sach
            score = 1 / (1 + new_pts)             -- luon giai quyet 1 vi pham moi
            if score > best_score:
                best_score = score
                best_idx   = i

        if best_idx = None: break  -- khong con goi nao vua

        -- Chon goi best_idx
        p = packages[best_idx]
        sel_viol.add(p.violation)
        sel_points.add_all(p.points)
        total_pts += count new points added
        selected.add(best_idx)

    -- FALLBACK: dam bao luon co tien trien
    if selected is empty:
        oldest = package with smallest first_seen
        selected.add(index of oldest)

    return selected
\end{verbatim}

\textbf{Fallback bắt buộc:} Nếu $P_{\mathrm{sel}}$ vẫn rỗng sau vòng chọn lọc,
chương trình sẽ bị lặp vô hạn. Xử lý: lấy gói có \texttt{first\_seen} nhỏ nhất
và thêm bắt buộc dù có vượt ngân sách.

Giá trị $K$ để thử: $8, 32, 128$ và thêm tất cả ($\infty$).
Kiểm thử với $K=\infty$ trước để đảm bảo kết quả giống baseline.

\subsection{Cách mốc được biểu diễn trong solver}
\label{sec:encoding}

Với các mốc $\tau_1 < \tau_2$ trong $\Lambda_v$, biến $d_v(\tau)$ mang nghĩa
$t_v \ge \tau$ và phải thỏa:
\begin{equation}
 d_v(\tau_2) \Rightarrow d_v(\tau_1)
 \quad\text{(clause: }\neg d_v(\tau_2) \lor d_v(\tau_1)\text{).}
\label{eq:ladder}
\end{equation}

Nếu lịch vi phạm $t_b \ge t_a + h$ và đang chọn $t_a = \alpha$, gói cơ bản:
\begin{align}
 \Lambda_b &\leftarrow \Lambda_b \cup \{\alpha+h\},\\
 &\neg d_a(\alpha) \lor d_b(\alpha+h).
\end{align}

Ví dụ $t_a=10$, $h=5$: thêm mốc 15 cho $b$ và implication
$d_a(10) \Rightarrow d_b(15)$. Lịch cũ không thể giữ nguyên vì vế trái đúng,
vế phải sai.

\textit{Kết quả cần có:} Chính sách bật/tắt được; không lặp vô hạn; log số
vấn đề, số gói đề xuất, số gói chọn, số mốc thực sự thêm mỗi vòng.

% =============================================================================
\section{Cải tiến 3: chia sẻ mã hóa AMO}
\label{sec:improve3}
% =============================================================================

\subsection{Đây là hướng khảo sát có điều kiện}

Tại mỗi thời điểm, mỗi đoạn đường chỉ cho tối đa 1 tàu đi (AMO --- At Most One).
Ở các thời điểm gần nhau, danh sách tàu tranh nhau thường gần giống nhau. Nếu
hai nhóm có phần chung $S$ lớn, có thể tạo biến chung $g$ thay vì mã hóa lại
từ đầu.

Tuy nhiên, chia sẻ sai có thể loại bỏ lịch hợp lệ. Vì vậy sinh viên phải đo
dữ liệu thực tế trước, rồi mới cài đặt.

\subsection{Bước khảo sát cần làm trước}
\begin{enumerate}
\item Xuất danh sách tất cả nhóm AMO trong quá trình chạy.
\item Đo mức độ giao nhau giữa các nhóm liên tiếp.
\item Ước lượng số biến và clause có thể tiết kiệm.
\item Chỉ tiếp tục nếu: nhiều nhóm có giao nhau đáng kể \emph{và} chi phí
      tạo cấu trúc dùng chung thấp hơn lợi ích.
\end{enumerate}

\subsection{Hai cách mã hóa cần có để so sánh}

\textbf{Pairwise} ($n < \theta$, không cần biến phụ):
\[
 \neg x_i \lor \neg x_j\quad(1 \le i < j \le n),
 \qquad n(n-1)/2\text{ clause.}
\]

\textbf{Sequential Counter} ($n \ge \theta$, tạo $3n-4$ clause, $n-1$ biến phụ):
\begin{align}
 &\neg x_1 \lor s_1,\label{eq:sc}\\
 &\neg x_i \lor s_i,\;\;
   \neg s_{i-1} \lor s_i,\;\;
   \neg x_i \lor \neg s_{i-1}
   \quad(2 \le i \le n-1),\notag\\
 &\neg x_n \lor \neg s_{n-1}.\notag
\end{align}

Lưu ý: code hiện tại có thêm chiều đảo của prefix nên thực tế tạo $4n-5$ clause.
Đây là biến thể khác, phải ghi rõ trong cấu hình khi so sánh.

\subsection{Chia sẻ phần chung $S$}

Nếu $X = A \cup S$ và $Y = S \cup B$ ($A, S, B$ rời nhau), tạo biến
$g \Leftrightarrow \bigvee_{s \in S} s$, rồi thay bằng:
\[
 \AMO(S) \;\land\; \AMO(A \cup \{g\}) \;\land\; \AMO(B \cup \{g\}).
\]
Nếu $S = \emptyset$: không chia sẻ. Nếu $|S|=1$: dùng phần tử đó thay $g$.

Chỉ đưa vào cùng AMO các indicator đã xác nhận đôi một xung đột. Hai biến
của cùng hoạt động ở hai thời điểm khác nhau là hai biến khác nhau.

\textit{Kết quả cần có:} Trước hết là báo cáo mức độ giao nhau. Prototype chỉ
làm nếu báo cáo cho thấy có triển vọng.

% =============================================================================
\section{Tích hợp và kiểm thử}
\label{sec:integrate}
% =============================================================================

\subsection{Luồng chương trình sau khi tích hợp}

Sau khi hoàn thành riêng từng phần, luồng chính là:

\begin{center}
\textit{Đọc dữ liệu}
$\to$ \textit{Suy luận sớm (Cải tiến 1)}
$\to$ \textit{Giải MaxSAT}
$\to$ \textit{Kiểm tra lịch + chi phí}
$\to$ \textit{Chọn refinement (Cải tiến 2)}
$\to$ quay lại suy luận sớm.
\end{center}

Chương trình chỉ báo tối ưu khi $\LB = \UB$. Nếu hết thời gian: trả lịch
tốt nhất cùng hai cận và trạng thái timeout.

\subsection{Nguyên tắc kiểm thử}

Sau mỗi cải tiến, chạy lại cùng instance. Nghiệm cuối phải y hệt baseline
về chi phí. Số vòng lặp có thể giảm, không được tăng.

Viết một chương trình brute-force (thử tất cả lịch có thể) cho 2--3 tàu.
Chương trình này không dùng lại code của MaxSAT-DDD và đóng vai trò đáp án
đối chiếu cho mọi cải tiến.

\subsection{Bộ bài kiểm thử tối thiểu}

\begingroup\small
\begin{longtable}{@{}p{12mm}p{58mm}p{80mm}@{}}
\toprule
ID & Tình huống & Điều cần kiểm tra \\\midrule\endhead
T01 & Một tàu, 3 hoạt động & $L_v$ được truyền đúng dọc hành trình. \\
T02 & Hai thứ tự đều còn khả thi & Chương trình không tự ép thứ tự. \\
T03 & Chỉ một thứ tự khả thi & Phát hiện và ghi lý do bắt buộc. \\
T04 & Không thứ tự nào khả thi & Báo không khả thi, không treo. \\
T05 & Tàu phải chờ trên tài nguyên & Thời gian chiếm dụng được kiểm tra đủ. \\
T06 & Độ trễ nằm đúng sát ngưỡng & Chi phí Step/Round đúng ở hai phía ngưỡng. \\
T07 & Lịch hợp lệ nhưng chi phí chưa khớp & Chương trình tiếp tục refinement. \\
T08 & Giới hạn refinement rất nhỏ ($K=1$) & Vẫn có tiến triển nhờ fallback. \\
T09 & Cùng mốc đề xuất nhiều lần & Mốc không bị thêm lặp. \\
T10 & Hai AMO có phần chung & Chia sẻ không cấm lịch hợp lệ. \\
T11 & Hết thời gian trước khi tối ưu & Trả lịch tốt nhất + hai cận + timeout. \\
T12 & Dữ liệu bất thường & Dừng với thông báo đủ để debug. \\
\bottomrule
\end{longtable}
\endgroup

\subsection{Các cấu hình cần so sánh}

\begin{tabularx}{\textwidth}{@{}p{24mm}X@{}}
\toprule
Tên & Nội dung \\\midrule
C & Phiên bản hội nghị đã chốt. \\
C+B & Thêm suy luận sớm (Cải tiến 1). \\
C+R & Thêm refinement có chọn lọc (Cải tiến 2). \\
C+B+R & Kết hợp cả hai. \\
C+B+R+S & Thêm chia sẻ AMO (nếu hướng này được tiếp tục). \\
\bottomrule
\end{tabularx}

% =============================================================================
\section{Thiết kế thực nghiệm}
\label{sec:experiment}
% =============================================================================

\subsection{Chỉ số cần báo cáo}
\begin{itemize}
\item Số bài chứng minh được tối ưu; số bài tìm được lịch hợp lệ.
\item Thời gian chứng minh tối ưu: trung vị, phân tán, PAR-2.
\item Số vòng refinement và số mốc thêm mỗi vòng.
\item Thời gian dành cho: suy luận, tạo model, solver, kiểm tra nghiệm.
\end{itemize}

PAR-2 với giới hạn $T$ và $N$ bài:
\[
 \mathrm{PAR2} = \frac{1}{N}\sum_{i=1}^{N} p_i,\quad
 p_i = \begin{cases}
   t_i & \text{chứng minh tối ưu trong }T,\\
   2T  & \text{không chứng minh trong }T.
 \end{cases}
\]

Không dùng thời gian trung bình trên các bài đã giải để kết luận một phương
pháp tốt hơn: hai phương pháp có thể giải được hai tập bài khác nhau.

\subsection{Cách chạy công bằng}
\begin{itemize}
\item Chạy ít nhất 5 lần mỗi cấu hình; với bài dưới 100ms thì tăng lên 20 lần.
\item Xáo trộn thứ tự chạy để giảm ảnh hưởng cache.
\item Tham số chính sách chỉ được điều chỉnh trên bộ bài dành riêng cho tuning,
      không trên bộ test chính.
\item Ghi rõ: số luồng, seed, timeout, giới hạn RAM, phiên bản solver, phần cứng.
\end{itemize}

% =============================================================================
\section{Kế hoạch 10 tuần}
\label{sec:plan}
% =============================================================================

\begingroup\small
\begin{longtable}{@{}p{13mm}p{68mm}p{69mm}@{}}
\toprule
Tuần & Công việc chính & Kết quả cuối tuần \\\midrule\endhead
1 & Đọc 5 file ở Mục~\ref{sec:read}; chạy script audit.
  & In được số vòng lặp, số mốc/vòng, chi phí. \\
2 & Chốt baseline; đọc bài hội nghị và bài Croella.
  & Cấu hình nền chạy lại được; file ghi cấu hình. \\
3 & Xây bộ bài T01--T12; viết chương trình brute-force.
  & Có đáp án đúng để kiểm thử các cải tiến. \\
4 & Cài đặt \texttt{propagate\_bounds()} và \texttt{find\_forced\_orders()}.
  & Qua T01--T05; log số thứ tự bắt buộc tìm được. \\
5 & Thêm hard clause vào solver; thử \texttt{tighten\_ub\_from\_cost()}.
  & Nghiệm y hệt baseline; quyết định giữ Mức 2 không. \\
6 & Tách 3 pha phát hiện/lập kế hoạch/áp dụng trong refinement.
  & Với $K=\infty$: kết quả y hệt baseline. \\
7 & Cài đặt scoring, vòng chọn lọc và fallback.
  & Không lặp vô hạn; log đủ mỗi vòng. \\
8 & Tích hợp C+B+R; khảo sát mức độ giao nhau AMO.
  & Pilot ngắn; quyết định có làm Cải tiến 3 không. \\
9 & Chạy thí nghiệm chính trên toàn bộ benchmark.
  & Log đầy đủ; tách rõ tối ưu, timeout, lỗi. \\
10 & Phân tích kết quả; viết báo cáo.
  & Tái tạo được bảng/hình; xác định đóng góp thực tế. \\
\bottomrule
\end{longtable}
\endgroup

Nếu bước baseline hoặc kiểm thử phát hiện lỗi mô hình, ưu tiên sửa trước
khi tiếp tục tối ưu tốc độ.

% =============================================================================
\section{Báo cáo và bàn giao}
\label{sec:report}
% =============================================================================

\subsection{Sản phẩm cuối cùng}
\begin{itemize}
\item Mã nguồn có thể bật/tắt từng cải tiến qua tùy chọn dòng lệnh.
\item Bộ bài T01--T12 và chương trình kiểm tra brute-force.
\item File ghi cấu hình của từng chiến dịch thực nghiệm.
\item Dữ liệu thô và log theo từng vòng lặp.
\item Script sinh bảng và hình từ CSV.
\item Báo cáo chỉ rõ cải tiến nào có ích, trong trường hợp nào, giới hạn là gì.
\end{itemize}

\subsection{Mẫu báo cáo hằng tuần}

Mỗi báo cáo khoảng 1--2 trang, gồm: (1) mục tiêu tuần; (2) phần đã hoàn thành;
(3) cách kiểm tra và kết quả (kể cả kết quả âm); (4) vấn đề còn tồn tại;
(5) kế hoạch tuần tiếp theo. Kèm một lệnh chạy lại ví dụ quan trọng nhất.

\subsection{Tiêu chí hoàn thành}
\begin{enumerate}
\item Baseline đã xác minh; các sai khác cũ được giải thích.
\item Hai cải tiến chính bật/tắt độc lập và đều qua T01--T12.
\item Chương trình không báo tối ưu khi chưa có đủ bằng chứng ($\LB = \UB$).
\item Thí nghiệm có ablation, đo lặp và phản ánh cả timeout.
\item Báo cáo phân biệt rõ phần đã có ở bài hội nghị với phần mới.
\end{enumerate}

% =============================================================================
\section{Biên dịch tài liệu}
% =============================================================================

\begin{verbatim}
cd manuscript/journal
latexmk -xelatex student_implementation_guide_vi.tex
mkdir -p ../../output/pdf
cp student_implementation_guide_vi.pdf ../../output/pdf/
\end{verbatim}

Trên Overleaf: chọn compiler XeLaTeX.

\clearpage
\begin{thebibliography}{9}
\bibitem{croella}
A.~L. Croella, B.~Luteberget, C.~Mannino và P.~Ventura.
\emph{A MaxSAT approach for solving a new Dynamic Discretization Discovery
model for train rescheduling problems}.
Computers \& Operations Research, 167, 106679, 2024.
\url{https://doi.org/10.1016/j.cor.2024.106679}.

\bibitem{sinz}
C.~Sinz.
\emph{Towards an Optimal CNF Encoding of Boolean Cardinality Constraints}.
CP 2005.

\bibitem{source}
B.~Luteberget và cộng sự. Repository MaxSAT train scheduling.
\url{https://github.com/luteberget/maxsattrainscheduling}.

\bibitem{conference}
T.~V. Kieu, T.~H. Nguyen và K.~V. To.
\emph{An Efficient MaxSAT-DDD Approach for Train Rescheduling via
Precedence Propagation and Hybrid AMO Encodings}.
Bản thảo tại \pathref{manuscript/submission/main.tex}.
\end{thebibliography}
\end{document}
